\chapter{บทนำ}

\section{ความเป็นมาและความสำคัญของปัญหา}

ภาคการเกษตรมีบทบาทสำคัญในการขับเคลื่อนเศรษฐกิจไทย โดยมีแรงงานมากกว่าร้อยละ 30 ของกำลังแรงงานทั้งหมด อย่างไรก็ตาม ภาคการเกษตรกลับมีสัดส่วนเพียงร้อยละ 10 ของผลิตภัณฑ์มวลรวมในประเทศ (GDP) และมีอัตราการเติบโตที่ลดลงในช่วงหลายทศวรรษที่ผ่านมา

เกษตรกรไทยต้องเผชิญกับปัญหาหลายประการ ทั้งความไม่แน่นอนจากสภาพภูมิอากาศ การเสื่อมโทรมของทรัพยากรธรรมชาติ และข้อจำกัดในการเข้าถึงเทคโนโลยีและองค์ความรู้สมัยใหม่ ส่งผลให้ผลิตภาพต่ำและรายได้ลดลง

จากปัญหาดังกล่าว กระทรวงเกษตรและสหกรณ์จึงได้ริเริ่มโครงการ “เกษตรกรอัจฉริยะ” (Smart Farmer: SF) เพื่อยกระดับเกษตรกรให้สามารถพัฒนาความรู้ ทักษะ และการใช้เทคโนโลยีในการผลิตทางการเกษตรได้อย่างมีประสิทธิภาพ

อย่างไรก็ตาม จากการดำเนินโครงการที่ผ่านมา ยังพบว่าการจัดอบรมมีลักษณะเป็นแบบเดียวสำหรับเกษตรกรทุกกลุ่ม (One-size-fits-all) ซึ่งอาจไม่สอดคล้องกับบริบทและความต้องการเฉพาะของเกษตรกรแต่ละราย ส่งผลให้ประสิทธิผลของโครงการยังไม่เป็นไปตามเป้าหมาย

ดังนั้น งานวิจัยนี้จึงมุ่งเน้นการศึกษาผลกระทบของโครงการ Smart Farmer ที่มีต่อทักษะของเกษตรกร และศึกษาความเป็นไปได้ในการนำเทคโนโลยีการเรียนรู้ของเครื่อง (Machine Learning) มาใช้ในการจำแนกกลุ่มเกษตรกรเพื่อพัฒนาหลักสูตรการฝึกอบรมให้มีความเหมาะสมมากยิ่งขึ้น

\section{วัตถุประสงค์ของการวิจัย}

\begin{enumerate}
\item เพื่อศึกษาผลกระทบของการเข้าร่วมโครงการ Smart Farmer ที่มีต่อทักษะของเกษตรกร
\item เพื่อวิเคราะห์ความสัมพันธ์ระหว่างคุณลักษณะของหลักสูตรฝึกอบรมกับการพัฒนาทักษะของเกษตรกร
\item เพื่อพัฒนาแนวทางการจำแนกกลุ่มเกษตรกรโดยใช้เทคโนโลยี Machine Learning
\end{enumerate}

---

\section{ขอบเขตการวิจัย}

\textbf{ขอบเขตด้านประชากร}  
การวิจัยนี้มุ่งศึกษาเกษตรกรที่เข้าร่วมโครงการ Smart Farmer และเกษตรกรที่ไม่ได้เข้าร่วมโครงการ ในช่วงปี พ.ศ. 2565–2566

\textbf{ขอบเขตด้านกลุ่มตัวอย่าง}  
กลุ่มตัวอย่างในการวิจัยมีจำนวนประมาณ 350–400 ราย โดยใช้วิธีการจับคู่ด้วยคะแนนความโน้มเอียง (Propensity Score Matching)

\textbf{ขอบเขตด้านเนื้อหา}  
การศึกษานี้มุ่งเน้นการวิเคราะห์ทักษะของเกษตรกรก่อนและหลังการเข้าร่วมโครงการ และประเมินผลกระทบของโครงการ โดยใช้แบบจำลองเชิงเศรษฐมิติและเทคนิคการเรียนรู้ของเครื่อง

\section{กรอบแนวคิดในการวิจัย}

งานวิจัยนี้กำหนดให้  
\begin{itemize}
\item ตัวแปรตาม คือ ระดับทักษะของเกษตรกร (Skill)
\item ตัวแปรอิสระ ได้แก่ การเข้าร่วมโครงการ (Training) และคุณลักษณะส่วนบุคคลของเกษตรกร
\end{itemize}

โดยใช้แบบจำลอง Difference-in-Differences (DID) เพื่อประเมินผลกระทบของโครงการ

\section{คำจำกัดความที่ใช้ในการวิจัย}

\textbf{เกษตรกรอัจฉริยะ (Smart Farmer)}  
หมายถึง เกษตรกรที่มีความสามารถในการประยุกต์ใช้เทคโนโลยีและองค์ความรู้ในการผลิตทางการเกษตร

\textbf{การเรียนรู้ของเครื่อง (Machine Learning)}  
หมายถึง เทคนิคทางปัญญาประดิษฐ์ที่ใช้ในการวิเคราะห์ข้อมูลและจำแนกกลุ่มตัวอย่าง

\section{ประโยชน์ที่คาดว่าจะได้รับ}

\begin{enumerate}
\item ได้ข้อมูลเชิงประจักษ์เกี่ยวกับผลกระทบของโครงการ Smart Farmer
\item ได้แนวทางการพัฒนาหลักสูตรฝึกอบรมที่มีความเหมาะสมกับเกษตรกรแต่ละกลุ่ม
\item ได้โมเดล Machine Learning สำหรับการจำแนกกลุ่มเกษตรกร
\end{enumerate}